%% =====================================================================
%%  T-01-02  Three Different Activities
%%  -------------------------------------------------------------------
%%  One idea per file. Core is mandatory. Slots in canonical order.
%%  Max 700 words. No layout commands. No filler. New source material
%%  goes in sources/<BOOK>/T-01-02.tex -- never by rewriting this file.
%% =====================================================================
%% @kind: topic
%% @id: T-01-02
%% @title: Three Different Activities
%% @subtitle: Persuasion, manipulation and coercion are not the same move
%% @chapter: c01
%% @order: 20
%% @status: stable
%% @sources: B1, B3
%% @updated: 2026-09-19

\topic{T-01-02}{Three Different Activities}{Persuasion, manipulation and coercion are not the same move}

\begin{core}
Persuasion changes what a person wants. Manipulation changes what a person believes about the situation, so that what they want changes on its own. Coercion leaves what they want intact and makes it expensive. Confusing the three produces bad defences, because each is answered differently.
\end{core}
\begin{mechanism}
The distinction lies in what is being worked on and what the target is left holding. A persuaded person can restate your argument and still disagree; nothing was hidden. A manipulated person acts on a picture of the situation that you edited --- the alternatives, the urgency, the cost, who benefits --- and acts on it willingly, which is why they do not resist and why they often defend the decision afterwards. A coerced person complies while disagreeing, and retains the grievance. That is the practical difference: manipulation buys consent, coercion buys compliance, and consent is far more durable than compliance because it does not need to be re-applied.
\end{mechanism}
\begin{tells}
  \item Persuasion: the argument is on the table and you can repeat it back.
  \item Manipulation: the options, the timing or the stakes have been arranged, and you feel you chose freely.
  \item Coercion: the cost of refusing is stated or obvious, and the offer does not improve when you hesitate.
\end{tells}
\begin{moves}
  \item Name which of the three is in use before responding to it.
  \item Against persuasion, answer the argument.
  \item Against manipulation, reconstruct the situation independently: get the real options, the real deadline, the real price, from a source the other party does not control.
  \item Against coercion, stop arguing about the merits and work on the cost of compliance --- reduce your exposure, or make the cost of enforcing it higher.
\end{moves}
\begin{counters}
  \item Reconstruct the choice set from outside. Manipulation depends on you accepting the frame; a second source breaks it.
  \item Slow the clock. Manufactured urgency is the most common signal that the picture has been edited.
  \item Ask for the terms in a form you can keep. Consent built on an edited picture rarely survives being written down.
  \item For coercion, do not negotiate the framing at all --- negotiate the exit.
\end{counters}
\begin{limits}
The boundaries blur in practice, and both sources treat them as blurring. Selective honesty used to disarm someone is manipulation that looks like persuasion; a deadline can be real and manufactured at once. A taxonomy that is too clean makes you confident and wrong. Its use is diagnostic --- it tells you which question to ask next --- not verdictal.
\end{limits}

\sources{B1, B3}
\seesources
\seealso{T-01-01, T-14-01, T-19-01}
